\section{Spacelike exchange and locality}
\label{sec:locality}
We prove that exchanging two adjacent spacelike-separated insertions does
not change the boundary distribution. The argument keeps all the other
insertions present. It first makes the two tube functions analytic near a
partial boundary with one tied imaginary-time gap. Euclidean symmetry
identifies their germs when the real times vanish. Analytic continuation
along a path then identifies them at arbitrary spacelike configurations.
Finally their partial boundary distributions agree, and the remaining
imaginary heights can be removed.

Fix adjacent slots $i,j=i+1$ and let $\sigma$ exchange them. It acts on
tuples of labels, real points and complex points by the same permutation.
Put
\[
 O_{ij}=\{X:X_j-X_i\text{ is spacelike}\}
       =\{X:|X_j^0-X_i^0|<|\mathbf X_j-\mathbf X_i|\}.
\]
Choose scalar heights $y$ with $y_i=y_j$ and with every other ordered
pair strictly increasing. Then $d(y)\in C_n$ and $\sigma d(y)=d(y)$.
The point $X+i d(y)$ lies in the closures of both $\mathcal T_n$ and
$\sigma^{-1}\mathcal T_n$: strictly separate the tied heights in the
appropriate order by arbitrarily small amounts. The functions to be
compared are
\[
 z\longmapsto H_n(v;z),\qquad
 z\longmapsto H_n(v\circ\sigma;\sigma z).
\]

\subsection{Small imaginary boosts enter both tubes}
For a spatial unit vector $e$, write, in time-first coordinates,
\[
 C_e(t,\mathbf x)=(e\cdot\mathbf x,te),\qquad
 P_e(t,\mathbf x)=(t,(e\cdot\mathbf x)e).
\]
The boost and its entire complex extension are
$B_e(z)=I+\sinh z\,C_e+(\cosh z-1)P_e$.
For a spacelike vector $a$, the choice
$e=\mathbf a/|\mathbf a|$ gives
$C_ea=(|\mathbf a|,a^0e)\in V_+$.

\begin{lemma}\label{lem:imaginary-boost-path}
For $X\in O_{ij}$ and $q=X+i d(y)$ there are $e$ and $\tau>0$ such that
$B_e(i\tau)q\in\mathcal T_n$ and
$\sigma B_e(-i\tau)q\in\mathcal T_n$.
Moreover, for all $z$ sufficiently close to $q$ in $\mathcal T_n$,
$B_e(ib)z\in\mathcal T_n$ for every $0\leq b\leq\tau$.
For all $z$ sufficiently close to $q$ in $\sigma^{-1}\mathcal T_n$,
$\sigma B_e(-ib)z\in\mathcal T_n$ throughout the same interval.
\end{lemma}
\begin{proof}
Take $e$ from the spacelike difference $X_j-X_i$ as above.
For any adjacent difference $z_{k+1}-z_k$, write its real and imaginary
parts as $x$ and $Y$. The imaginary part after the positive imaginary
boost is
\begin{equation}\label{eq:boost-imaginary-gap}
 Y+\sin b\,C_ex+(\cos b-1)P_eY.
\end{equation}
At $q$, every gap other than $(i,j)$ has $Y\in V_+$. By openness and
continuity those finitely many transformed gaps remain in $V_+$ for
$(b,z)$ in a product neighborhood of $(0,q)$.
At the distinguished gap $Y=0$ and $C_ex\in V_+$. For small $b$ rewrite
\eqref{eq:boost-imaginary-gap} as
\[
 Y+\sin b\left(C_ex-\frac{\sin b}{1+\cos b}P_eY\right).
\]
The bracket is continuous near $(0,q)$ and lies in $V_+$ there, after
shrinking the neighborhood. For $0<b<\pi$, the sine is positive.
If $Y\in\overline V_+$, the sum is strictly in $V_+$: the spatial
triangle inequality shows that a closed forward-cone vector plus a
positive multiple of a strict one is strict.

For the negative boost followed by the exchange, the original tied gap
is reversed and the sine changes sign. Its inward vector is again
$C_e(X_j-X_i)$. All other imaginary gaps at $q$ remain strictly forward,
since the tied heights agree and are the only tie. The same neighborhood
argument applies. Take a common product neighborhood for both orientations
and all gaps, and choose $\tau>0$ inside its positive parameter interval.
At $q$ the closed-gap condition gives the two strict endpoints. At nearby
interior points it gives strictness for $b>0$, while $b=0$ follows from
membership in the relevant original tube. This proves the asserted
uniform path property, including its endpoint at zero.
\end{proof}

\subsection{Holomorphic germs across the tied gap}
The labels in any finite tuple lie in one common finite-dimensional level.
On that level the continuous representation $r\mapsto\rho(B_e(r))$
has an entire extension. Here is the elementary argument, which uses
operator-norm continuity only in this finite-dimensional space.
For a continuous one-parameter group $R(r)$ in its matrix algebra, put
$A(t)=\int_0^tR(s)\,ds$. Then $A(t)/t\to I$ as $t\to0$, so $A(a)$
is invertible for a sufficiently small nonzero $a$. The group law gives
\[
 R(t)A(a)=A(t+a)-A(t).
\]
The right side is continuously differentiable. Thus $R$ is differentiable;
differentiating the group law gives $R'(t)=R(t)K$, $K=R'(0)$.
The derivative of $R(t)e^{-tK}$ is zero and its value at zero is $I$.
Therefore $R(t)=e^{tK}$, and $R(z)=e^{zK}$ is the required entire extension.
This construction is compatible with the finite-level inclusions, since
the real representations are compatible and the identity theorem extends
their equality to complex parameters.

For a point $q$ and parameters supplied by
Lemma~\ref{lem:imaginary-boost-path}, define near $q$
\begin{align*}
 A_q(z)&=H_n(R(i\tau)v;B_e(i\tau)z),\\
 B_q(z)&=H_n(R(-i\tau)(v\circ\sigma);\sigma B_e(-i\tau)z).
\end{align*}
They are holomorphic in a neighborhood of $q$, because their arguments
at $q$ lie strictly in the tube. They extend the respective original
tube functions. To prove this for $A_q$, fix a nearby
$z\in\mathcal T_n$ for which the entire path stays in the tube.
The function of one complex parameter
\[
 \zeta\longmapsto H_n(R(\zeta)v;B_e(\zeta)z)
\]
is holomorphic on an open neighborhood of the segment $i[0,\tau]$.
Moving labels cause no difficulty: expansion in a basis of the common
finite level makes the expression a finite sum of holomorphic coefficients
times the fixed-label holomorphic tube functions. For real $\zeta$ near
zero the expression equals $H_n(v;z)$ by tube covariance. The identity
theorem on the connected component containing the segment makes it constant
along that segment. Its value at $i\tau$ proves the assertion.
The same proof with the negative segment and the exchange gives the
assertion for $B_q$.

Such a germ is uniquely determined by its original tube function. Indeed,
two candidate germs are holomorphic on a common ball about $q$ and agree
on the intersection of that ball with the relevant open tube. This
intersection contains a nonempty open set because $q$ is in the tube's
closure. The identity theorem on the ball makes the germs equal.

\begin{lemma}[Equality on the Euclidean slice]\label{lem:euclidean-germs}
If every real time coordinate of $X$ is zero, then the germs $A_q$ and
$B_q$ coincide at $q=X+i d(y)$.
\end{lemma}
\begin{proof}
Let $M$ be the real-linear Wick map that sends Euclidean time to imaginary
time and leaves spatial coordinates real. Put
$p=((y_k,\mathbf X_k))_{k=1}^n$, so $Mp=q$.
This is a distinct-point Euclidean configuration. All pairs other than
$(i,j)$ have different time coordinates, and the distinguished pair has
different spatial coordinates because its real separation is spacelike.
The Euclidean correlation $S_v$ is consequently smooth near $p$.

On the open side where the Euclidean times are strictly ordered,
$A_q\circ M=S_v$ near $p$, by the tube restriction and pure-imaginary
identification. On the open side where their exchanged tuple is strictly
ordered, $B_q\circ M=S_v$, using Euclidean permutation symmetry as well.
The point $p$ is in the closure of each side. Differentiate the respective
identities on their open sides. For each order, continuity of the
derivatives up to $p$ shows that the two real derivatives at $p$ both
equal the corresponding derivative of $S_v$.

The image of $M$ spans the full complex coordinate space over $\mathbb C$:
each spatial basis vector is already in the image, and each time basis
vector is $-i$ times an image vector. By the chain rule the real derivatives
just compared are the restrictions of the complex multilinear derivatives
of $A_q$ and $B_q$ to this image. Complex multilinearity and spanning make
all their full complex derivatives equal. Their convergent local power
series therefore agree. This proves equality of germs, rather than merely
equality of values on the Euclidean slice.
\end{proof}

The equality propagates to every $X\in O_{ij}$. Keep its spatial coordinates
fixed and multiply all its real time coordinates by $s\in[0,1]$.
The distinguished separation remains spacelike throughout. Along
$q(s)=((sX_k^0,\mathbf X_k))_k+i d(y)$ both germs exist, and both
underlying tubes have $q(s)$ in their closures.
Germ equality is locally constant in $s$. To verify this, choose a common
ball where the two germs at $q(s_0)$ are holomorphic and extend their
original tube functions. For $s$ near $s_0$, uniqueness of germs identifies
the germs at $q(s)$ with these same functions. The identity theorem on
the ball says that equality near either point is equivalent to equality
near the other. Thus the set of parameters with equal germs and its
complement are both open in $[0,1]$. Connectedness and
Lemma~\ref{lem:euclidean-germs} at $s=0$ give equality at $s=1$.

\subsection{From germs to boundary distributions}
Let $W_{d(y),v}$ denote the partial boundary distribution
\eqref{eq:partial-damping}. Near a given $X\in O_{ij}$ it is represented
by the germ $A_q$ at the partial height:
\begin{equation}\label{eq:germ-boundary-integral}
 W_{d(y),v}(f)=\int A_q(Z+i d(y))f(Z)\,dZ
\end{equation}
for compactly supported smooth $f$ in a sufficiently small ball about $X$.
Indeed, choose a smaller product neighborhood of $(X,d(y))$ on which
the germ is bounded and agrees with the tube function at interior heights.
For $Y\to d(y)$ through $C_n^\circ$, its integrals against $f$ converge
to the displayed integral by dominated convergence, using a constant times
$|f|$ as majorant. The tube boundary theorem identifies the same limit
with $W_{d(y),v}(f)$. Applying this argument to the exchanged tube, and
changing variables by the coordinate permutation of absolute Jacobian one,
gives the analogous formula for
$W_{\sigma d(y),v\circ\sigma}(f\circ\sigma)$ with germ $B_q$.
Since these germs agree, the two distributions agree locally on $O_{ij}$.

They agree on every Schwartz test whose closed support is contained in
$O_{ij}$. For compact support, cover it by finitely many of the local
balls and apply a smooth partition of unity. For arbitrary Schwartz
support contained in this open set, multiply by compact cutoffs approaching
one in the Schwartz topology; their supports remain in $O_{ij}$ and
continuity passes the equality to the limit.

\begin{theorem}[Adjacent spacelike exchange]\label{thm:spacelike-exchange}
For every joint Schwartz test $f$ supported in $O_{ij}$,
\[
 W_n(v)(f)=W_n(v\circ\sigma)(f\circ\sigma).
\]
Consequently the reconstructed field operators commute on their common
domain whenever their test functions have mutually spacelike supports.
\end{theorem}
\begin{proof}
There is a sequence of tied profiles of the kind used above tending to zero.
For example, with slots numbered $1,\ldots,n$, take
$y_k^{(m)}=a_k/(m+1)$ where $a_k=k$ except $a_j=i$.
All other ordered pairs have positive gaps, and the exchanged profile is
unchanged. The preceding partial-distribution equality holds for each $m$.
Continuity on the closed shift set, and $W_0=W_n(v)$, give the asserted
equality by letting $m\to\infty$.

For the operator conclusion take any two exterior words and place the
two test functions $f_i,f_j$ between them. The closed support of their
joint product test is contained in the inverse image of
$\operatorname{supp}f_i\times\operatorname{supp}f_j$ under the two
coordinate projections: wherever either factor vanishes on a neighborhood,
so does the product. Mutually spacelike supports therefore put that joint
support in $O_{ij}$. The distributional identity exchanges exactly these
two letters and retains all spectators. Proposition~\ref{prop:word-locality}
now proves that the commutator vanishes on $\mathcal D$.
\end{proof}
